\documentclass[10pt,conference]{IEEEtran}
\IEEEoverridecommandlockouts
% The preceding line is only needed to identify funding in the first footnote. If that is unneeded, please comment it out.
%Template version as of 6/27/2024

\usepackage{svg}
\usepackage{cite}
\usepackage{amsmath,amssymb,amsfonts}
\usepackage{algorithmic}
\usepackage{graphicx}
\usepackage{textcomp}
\usepackage{xcolor}
\usepackage{tabularx}
\usepackage{float}

\def\BibTeX{{\rm B\kern-.05em{\sc i\kern-.025em b}\kern-.08em
    T\kern-.1667em\lower.7ex\hbox{E}\kern-.125emX}}
\begin{document}

\title{Analyzing the historical correlation of GitHub repositories features and popularity}

\author{\IEEEauthorblockN{1\textsuperscript{st} Felipe Ribeiro de Souza}
\IEEEauthorblockA{\textit{Universidade de São Paulo (IME-USP)}\\
São Paulo, Brazil \\
felipers@usp.br}
\and
\IEEEauthorblockN{2\textsuperscript{nd} Higor Amário de Souza}
\IEEEauthorblockA{\textit{Universidade de São Paulo (POLI-USP)}\\
São Paulo, Brazil \\
higoramario@usp.br}
}

\maketitle

\begin{abstract}
TODO
\end{abstract}

\begin{IEEEkeywords}
T, O, D, O.
\end{IEEEkeywords}

\section{Introduction}
TODO

\section{Literature Review}
Studies based on repository features can be very different from each other, specially when it comes to tools like GitHub, where a large number of functionalities can be found. In this section we will go through the state of art and present how our work can complement it.

\subsection{Related Work}
Features like stars, contributors, forks, commits and project age have been studied in \cite{factors.impact.popularity}. It points that no correlation was found between stars and repository age, a weak correlation was noticed between stars, commits and contributors, and a strong correlation between stars and forks.

The correlation of popularity metrics such as stars and forks has been studied in \cite{readme.correlation}, the same study also discuss the correlation inside README.md file atributes - number of lines, number of images, update interval, etc - and how it is related to the popularity of the repository. One of the study conclusions is that the repository popularity is directly correlated with the frequency and number of updates in README.md files.

The correlation analysis presented by \cite{random.tree} is an interesting example of how correlation is helpful in telling which data should be discarded when training prediction models. Among of other features analyzed, developers, issues and pull requests is a point in common with our work, the difference is that we look at the authors who performed the issues and pull requests actions, and not just at the events themselves.

Alongside with other 92 factors, \cite{pull.request.decisions} studied how developer and integrator relate to each other, and how this relations can influence on the acceptance rate of pull requests.

\subsection{Contributions}
Software repositories features - such as issues, pull requests, popularity, commits, among others - have already been extensively studied in different contexts. However, to the best of our knowledge, no study has yet been published that correlates issues and pull requests (as well as their respective authors and reviewers) with other repository popularity metrics.

\section{Methodology}
This section presents the process of data acquisition, preparation, synthesis and analysis, according to the best practices for Mine Software Repositories (MSR) proposed by \cite{cookbook}.

\subsection{Data Acquisition and Preparation}
\label{acquisition}
Prior to the acquisition phase, we first needed to define a systematic and reproducible method for selecting which GitHub repositories would be the subject of our work.

We used the SEART GitHub search tool \cite{seart} to find the repositories that best fit our purposes. The following filters where applied: Non-fork repositories created before 1\textsuperscript{st} January 2021, with active issues and pull requests functionalities, containing at least 1000 commits, 500 issues, 500 pull requests, 500 stars, 100 forks and 1 release. Resulting in a sample of 6448 repositories. Then we ranked the top 10 main languages of the sample, which were: Python (16.10\% of the repositories), TypeScript (13.62\%), Go (10.72\%), C++ (10.02\%), JavaScript (9.65\%), Java (8.82\%), C\# (5.20\%), C (4.93\%), PHP (4.39\%) and Rust (4.03\%). As our objective for this work was to analyze popular and well-stablished projects, the repositories were ordered by the number of stars, and the top 15 repositories of each language were selected. To perform the final selection, we analyzed all the 150 resulting repositories and picked 3 library or framework projects of each language.

With our 30 projects in hand, the next step was to collect the history for each one of the repository features. To achieve this, we rely on the use of the follow tools:

\begin{itemize}
    \item PyDriller \cite{pydriller}, a Python library used to collect data available within Git itself, does not requires communication with the GitHub API.
    \item PoolinGH \cite{poolingh} is a JavaScript library used to collect data that requires communication with the GitHub API. To make better use of the rate-limit consumption, we made use from both Rest and GraphQL APIs. This library, in particular, was very useful to queue the thousands of requests we had to make and handle the errors returned by GitHub APIs.
    \item pytrends-modern\footnote{https://github.com/yiromo/pytrends-modern} is a Python library used to collect the Google Trends popularity data for the specified repository. This tool is a sucessor for the archived pytrends\footnote{https://github.com/GeneralMills/pytrends} library.
\end{itemize}

In addition to these tools, we also used the Python libraries Pandas\footnote{https://github.com/pandas-dev/pandas}, NumPy\footnote{https://github.com/numpy/numpy}, Matplotlib\footnote{https://github.com/matplotlib/matplotlib} and SciPy\footnote{https://github.com/scipy/scipy} to support the work presented in the sections \ref{synthesis} and \ref{analysis}.

Table \ref{table:features} details all collected data, their respective source and mining tool used.

% The follow table was generated in https://www.tablesgenerator.com, use \begin{table}[H] to pin the table in this part of the text
\begin{table}[H]
\caption{GitHub Repository Features Collected}
\label{table:features}
\centering
\begin{tabularx}{\linewidth}{|>{\centering\arraybackslash}X|c|c|}
\hline
\textit{\textbf{Feature}}          & \textit{\textbf{Source}} & \textit{\textbf{Mining Tool}} \\ \hline
Commits                            & GitHub                  & PyDriller                     \\ \hline
Commit Authors                     & GitHub                  & PyDriller                     \\ \hline
Forks                              & GitHub                  & PoolinGH                      \\ \hline
Created Issues                     & GitHub                  & PoolinGH                      \\ \hline
Closed Issues                      & GitHub                  & PoolinGH                      \\ \hline
Authors that Created Issues        & GitHub                  & PoolinGH                      \\ \hline
Authors that Closed Issues         & GitHub                  & PoolinGH                      \\ \hline
Created Pull Requests              & GitHub                  & PoolinGH                      \\ \hline
Closed Pull Requests               & GitHub                  & PoolinGH                      \\ \hline
Merged Pull Requests               & GitHub                  & PoolinGH                      \\ \hline
Authors that Created Pull Requests & GitHub                  & PoolinGH                      \\ \hline
Authors thats Closed Pull Requests & GitHub                  & PoolinGH                      \\ \hline
Authors that Merged Pull Requests  & GitHub                  & PoolinGH                      \\ \hline
Stars                              & GitHub                  & PoolinGH                      \\ \hline
Relative Interest                  & Google                  & pytrends-modern               \\ \hline
\end{tabularx}
\end{table}

It is worth noting that the inclusion of watchers would also be an interesting feature for this study, however, GitHub does not provide a way to collect the historical data regarding this feature. This means that it is only possible to obtain the number of watchers as it stands at the time of data collection, making our analysis impracticable for this feature.

With the historical data collected for each repository, a crucial step before continuing the analysis is the normalization of the dataset. First of all, we created a monthly view for each one of the features, aggregating the authors and summating the events that was daily distributed. If a feature had no events recorded for a given month, we consider the value for that month to be zero. Next, we merged all the different features into an unified view for each project, to make further analysis simpler.

\subsection{Synthesis}
\label{synthesis}
Before calculating the correlation of the data, we performed an Anderson-Darling (AD) test \cite{anderson.darling} for each feature in each repository. AD test is used to determine whether a dataset follows a specific distribution, in our case, we employed it to verify whether the data follows a normal distribution. We consider the distribution to be normal when the p-value is greater than 0.05, otherwise, we consider it non-normal.

With this information we applied Pearson's and Spearman's correlations to the data. For the pairs which both features represented normal distributions according to AD test, we considered Pearson's correlation, as it is more indicated for this use case. Otherwise, if none or only one of the features presents a normal distribution, Spearman's rank was considered in the final analysis.

\subsection{Analysis}
\label{analysis}
TODO

\section{Results}
TODO

\section{Future Work}
TODO

\section{Conclusions}
TODO

\begin{figure*}[t]
    \centering
    \includesvg[width=0.8\textwidth]{fig1.svg}
    \caption{Correlation analysis of GitHub repositories features.}
    \label{fig}
\end{figure*}

\bibliographystyle{plain}
\bibliography{references}

\end{document}
